\documentclass[11pt, a4paper]{article}
\usepackage{amsmath, amssymb, amsfonts}
\usepackage{geometry}
\usepackage{hyperref}
\usepackage{booktabs}
\usepackage{parskip}

\geometry{margin=1in}

% ---------- notation shortcuts ----------
\newcommand{\vect}[1]{\boldsymbol{#1}}
\newcommand{\qv}{\vect{q}_v}
\newcommand{\w}{\vect{\omega}}
\newcommand{\sgn}{\operatorname{sgn}}
\newcommand{\clip}{\operatorname{clip}}
\newcommand{\diag}{\operatorname{diag}}

\title{\textbf{Quaternion Attitude Dynamics, Nonlinear PD/PID Attitude Control \& Smooth Trajectory Tracking for a Rigid Body and Missile Pitch Channel}}
\author{}
\date{}

\begin{document}

\maketitle

\section{Attitude Representation with Unit Quaternions}

\subsection{Axis--Angle to Quaternion}
The attitude is stored as a scalar-first unit quaternion $\vect{q} = [q_0,\ q_1,\ q_2,\ q_3]^T = [q_0,\ \qv^T]^T$. For a rotation of angle $\phi$ about an axis $\vect{e}$ (normalised to unit length, $\hat{\vect{e}} = \vect{e}/\lVert\vect{e}\rVert$):
\begin{equation}
\vect{q} = \begin{bmatrix} \cos(\phi/2) \\ \hat{\vect{e}}\,\sin(\phi/2) \end{bmatrix},
\qquad
\vect{q}^{*} = \begin{bmatrix} \cos(\phi/2) \\ -\hat{\vect{e}}\,\sin(\phi/2) \end{bmatrix}
\end{equation}
where $\vect{q}^{*}$ is the conjugate (the inverse rotation for a unit quaternion). A zero-length axis ($\lVert\vect{e}\rVert < 10^{-12}$) is rejected.

\subsection{Quaternion (Hamilton) Product}
For $\vect{p} = [p_0,\ \vect{p}_v^T]^T$ and $\vect{q} = [q_0,\ \qv^T]^T$:
\begin{equation}
\vect{p} \otimes \vect{q} =
\begin{bmatrix}
p_0 q_0 - \vect{p}_v \cdot \qv \\
p_0\,\qv + q_0\,\vect{p}_v + \vect{p}_v \times \qv
\end{bmatrix}
\end{equation}

\subsection{Rotation of a Point}
A point $\vect{r} \in \mathbb{R}^3$ is embedded as the pure quaternion $\vect{p} = [0,\ \vect{r}^T]^T$ and rotated by the sandwich product; the rotated point is the vector part of the result:
\begin{equation}
\vect{p}' = \vect{q} \otimes \vect{p} \otimes \vect{q}^{*}
\end{equation}

\subsection{Error Quaternion}
With current attitude $\vect{q}$ and desired attitude $\vect{q}_d$ (both normalised first), the attitude error is the rotation from the desired frame to the body frame, $\vect{q}_e = \vect{q}_d^{*} \otimes \vect{q}$:
\begin{equation}
q_{e0} = q_0\, q_{d0} + \qv \cdot \vect{q}_{dv},
\qquad
\vect{q}_{ev} = q_{d0}\,\qv - q_0\,\vect{q}_{dv} + \qv \times \vect{q}_{dv}
\end{equation}
Perfect alignment corresponds to $\vect{q}_e = [\pm 1, 0, 0, 0]^T$. The scalar attitude-error angle used in all plots is
\begin{equation}
\theta_{\text{err}} = 2 \arccos\!\left( \lvert q_{e0} \rvert \right)
\end{equation}
The absolute value accounts for the double cover ($\vect{q}$ and $-\vect{q}$ are the same attitude), so $\theta_{\text{err}} \in [0, \pi]$.

\section{Secondary Attitude Representations}

\subsection{Direction Cosine Matrix}
The direction cosine matrix (inertial $\to$ body) built from the unit quaternion is
\begin{equation}
C(\vect{q}) =
\begin{bmatrix}
q_0^2 + q_1^2 - q_2^2 - q_3^2 & 2(q_1 q_2 + q_0 q_3) & 2(q_1 q_3 - q_0 q_2) \\
2(q_1 q_2 - q_0 q_3) & q_0^2 - q_1^2 + q_2^2 - q_3^2 & 2(q_2 q_3 + q_0 q_1) \\
2(q_1 q_3 + q_0 q_2) & 2(q_2 q_3 - q_0 q_1) & q_0^2 - q_1^2 - q_2^2 + q_3^2
\end{bmatrix}
\end{equation}

\subsection{Euler Angles (Yaw--Pitch--Roll, 3-2-1 Sequence)}
\begin{align}
\theta &= -\arcsin\!\big( 2 (q_1 q_3 - q_0 q_2) \big) && \text{(pitch, argument clipped to } [-1, 1]\text{)} \\
\phi &= \operatorname{atan2}\!\big( 2 (q_2 q_3 + q_0 q_1),\ q_0^2 - q_1^2 - q_2^2 + q_3^2 \big) && \text{(roll)} \\
\psi &= \operatorname{atan2}\!\big( 2 (q_1 q_2 + q_0 q_3),\ q_0^2 + q_1^2 - q_2^2 - q_3^2 \big) && \text{(yaw)}
\end{align}

\subsection{Modified Rodrigues Parameters}
The MRP vector is $\vect{\sigma} = \hat{\vect{e}} \tan(\phi/4)$. To stay away from the singularity at $\phi = 360^\circ$, the shadow set is selected whenever $q_0 < 0$:
\begin{equation}
\vect{\sigma} =
\begin{cases}
\dfrac{\qv}{1 + q_0}, & q_0 \ge 0 \\[2ex]
\dfrac{-\qv}{1 - q_0}, & q_0 < 0
\end{cases}
\end{equation}
which guarantees $\lVert \vect{\sigma} \rVert \le 1$.

\section{Rotational Equations of Motion}

\subsection{Quaternion Kinematic Differential Equation}
With body angular velocity $\w = [\omega_x, \omega_y, \omega_z]^T$, the kinematics $\dot{\vect{q}} = \tfrac{1}{2}\, \vect{q} \otimes [0,\ \w^T]^T$ are written in matrix form and augmented with a norm-constraint stabilisation term:
\begin{equation}
\dot{\vect{q}} = \frac{1}{2}\, \Omega(\w)\, \vect{q} \;+\; c \left( 1 - \vect{q}^T \vect{q} \right) \vect{q},
\qquad
\Omega(\w) =
\begin{bmatrix}
0 & -\omega_x & -\omega_y & -\omega_z \\
\omega_x & 0 & \omega_z & -\omega_y \\
\omega_y & -\omega_z & 0 & \omega_x \\
\omega_z & \omega_y & -\omega_x & 0
\end{bmatrix}
\end{equation}

\begin{itemize}
    \item $\tfrac{1}{2}\Omega(\w)\vect{q}$ --- \textbf{Kinematic Term}: Since $\Omega$ is skew-symmetric, this term alone preserves $\lVert \vect{q} \rVert$ exactly in continuous time.
    \item $c\,(1 - \vect{q}^T\vect{q})\,\vect{q}$ --- \textbf{Constraint (Baumgarte-type) Term}: Vanishes on the unit sphere and pulls the numerical solution back towards $\lVert \vect{q} \rVert = 1$ when integration error causes drift. The gain is $c = 1$ in all simulations.
\end{itemize}

\subsection{Euler's Rigid-Body Equation}
For inertia matrix $I$ and total body-frame torque $\vect{\tau}$:
\begin{equation}
I \dot{\w} = \vect{\tau} - \w \times (I \w)
\quad\Longrightarrow\quad
\dot{\w} = I^{-1} \big( \vect{\tau} - \w \times (I \w) \big)
\end{equation}
The term $\w \times (I\w)$ is the gyroscopic (cross-coupling) torque. The linear solve is performed either with a pre-computed inverse $I^{-1}$ or, optionally, with a pre-computed LU factorisation $P I = L U$ followed by forward/back substitution.

\subsection{Coupled Nonlinear State-Space Form}
Stacking attitude and rate gives the 7-dimensional state $\vect{x} = [\vect{q}^T,\ \w^T]^T \in \mathbb{R}^7$:
\begin{equation}
\dot{\vect{x}} = f(\vect{x}, \vect{\tau}) =
\begin{bmatrix}
\tfrac{1}{2}\, \Omega(\w)\, \vect{q} + c \left( 1 - \vect{q}^T \vect{q} \right) \vect{q} \\[1ex]
I^{-1} \big( \vect{\tau} - \w \times (I \w) \big)
\end{bmatrix}
\end{equation}

\subsection{Conservation-Law Checks}
For verifying the integrator in torque-free motion:
\begin{equation}
T_{\text{rot}} = \frac{1}{2}\, \w^T I\, \w,
\qquad
\vect{H}_B = I \w,
\qquad
\vect{H}_I = C(\vect{q})^T\, \vect{H}_B
\end{equation}
With $\vect{\tau} = \vect{0}$, the rotational kinetic energy $T_{\text{rot}}$ and the inertial angular momentum $\vect{H}_I$ must remain constant (the body-frame components $\vect{H}_B$ generally do not).

\section{Numerical Integration: Classical Fourth-Order Runge--Kutta}

\subsection{Kinematics-Only Integrator}
When $\w(t)$ is a prescribed function of time, only $\vect{q}$ is integrated, with step size $h$ and $g(\vect{q}, \w)$ denoting the right-hand side of the kinematic equation:
\begin{align}
\vect{k}_1 &= h\, g\big(\vect{q}_n,\ \w(t_n)\big) \\
\vect{k}_2 &= h\, g\big(\vect{q}_n + \tfrac{1}{2}\vect{k}_1,\ \w(t_n + \tfrac{h}{2})\big) \\
\vect{k}_3 &= h\, g\big(\vect{q}_n + \tfrac{1}{2}\vect{k}_2,\ \w(t_n + \tfrac{h}{2})\big) \\
\vect{k}_4 &= h\, g\big(\vect{q}_n + \vect{k}_3,\ \w(t_n + h)\big) \\
\tilde{\vect{q}}_{n+1} &= \vect{q}_n + \tfrac{1}{6} \left( \vect{k}_1 + 2\vect{k}_2 + 2\vect{k}_3 + \vect{k}_4 \right),
\qquad
\vect{q}_{n+1} = \frac{\tilde{\vect{q}}_{n+1}}{\lVert \tilde{\vect{q}}_{n+1} \rVert}
\end{align}

\subsection{Coupled-State Integrator}
For the full state $\vect{x}$ with torque law $\vect{\tau}(t, \vect{x})$:
\begin{align}
\vect{k}_1 &= f\big(\vect{x}_n,\ \vect{\tau}(t_n, \vect{x}_n)\big) \\
\vect{k}_2 &= f\big(\vect{x}_n + \tfrac{h}{2}\vect{k}_1,\ \vect{\tau}(t_n + \tfrac{h}{2},\ \vect{x}_n + \tfrac{h}{2}\vect{k}_1)\big) \\
\vect{k}_3 &= f\big(\vect{x}_n + \tfrac{h}{2}\vect{k}_2,\ \vect{\tau}(t_n + \tfrac{h}{2},\ \vect{x}_n + \tfrac{h}{2}\vect{k}_2)\big) \\
\vect{k}_4 &= f\big(\vect{x}_n + h\,\vect{k}_3,\ \vect{\tau}(t_n + h,\ \vect{x}_n + h\,\vect{k}_3)\big) \\
\vect{x}_{n+1} &= \vect{x}_n + \frac{h}{6} \left( \vect{k}_1 + 2\vect{k}_2 + 2\vect{k}_3 + \vect{k}_4 \right)
\end{align}
followed by re-normalisation of the quaternion part only, $\vect{q}_{n+1} \leftarrow \vect{q}_{n+1} / \lVert \vect{q}_{n+1} \rVert$.

\textbf{Torque handling inside a step.} Two modes are available:
\begin{itemize}
    \item \textbf{Continuous feedback} (default): the torque law is re-evaluated at every RK4 stage, as written above.
    \item \textbf{Zero-order hold}: the torque is evaluated once, $\vect{\tau}_n = \vect{\tau}(t_n, \vect{x}_n)$, and the same value is used in all four stages. This models a sampled-data (digital) controller.
\end{itemize}

\subsection{Propagation Over a Time Interval}
To propagate from $t_0$ to $t_f$ with $T = t_f - t_0 \ge 0$ and tolerance $\varepsilon = 10^{-9}$:
\begin{enumerate}
    \item Number of full steps: $n_{\text{full}} = \left\lfloor T/h + \varepsilon \right\rfloor$.
    \item Remainder: $r = T - n_{\text{full}}\, h$. A final partial step of size $r$ is taken only if $r > \varepsilon h$.
    \item Take $n_{\text{full}}$ steps of size $h$. After each step the time is reset to $t_i = t_0 + i\,h$ (rather than accumulated as $t_{i-1} + h$) to avoid floating-point round-off drift.
    \item If required, take the partial step and set the final time exactly to $t_f$.
    \item Return the time history and the $(n+1) \times 7$ state history, including the initial state.
\end{enumerate}

\section{Attitude Regulation: Quaternion-Feedback PD Law}

\subsection{Control Law}
For regulation to a fixed target attitude $\vect{q}_d$:
\begin{equation}
\vect{\tau} = -K_p\, \sgn(q_{e0})\, \vect{q}_{ev} \;-\; K_d\, \w,
\qquad
\sgn(q_{e0}) = \begin{cases} +1, & q_{e0} \ge 0 \\ -1, & q_{e0} < 0 \end{cases}
\end{equation}

\begin{itemize}
    \item $-K_p\, \sgn(q_{e0})\, \vect{q}_{ev}$ --- \textbf{Proportional (Attitude) Term}: Restoring torque along the error rotation axis. Since $\vect{q}_{ev} = \hat{\vect{e}} \sin(\theta_{\text{err}}/2)$, it behaves as a nonlinear spring.
    \item $\sgn(q_{e0})$ --- \textbf{Shortest-Path Selection}: $\vect{q}_e$ and $-\vect{q}_e$ represent the same attitude. The sign factor always drives the body through the rotation of less than $180^\circ$ and prevents unwinding.
    \item $-K_d\, \w$ --- \textbf{Derivative (Rate Damping) Term}: Dissipates rotational kinetic energy. $K_d$ may be a scalar or a $3 \times 3$ matrix.
\end{itemize}

\subsection{Lyapunov Stability Analysis}
The candidate function is the rotational kinetic energy plus an attitude ``potential'':
\begin{equation}
V(\vect{q}_e, \w) = \frac{1}{2}\, \w^T I\, \w + 2 K_p \left( 1 - \lvert q_{e0} \rvert \right)
\end{equation}
$V \ge 0$, with $V = 0$ only at $\w = \vect{0}$ and $q_{e0} = \pm 1$. For a constant target the error kinematics give $\dot{q}_{e0} = -\tfrac{1}{2}\, \vect{q}_{ev} \cdot \w$, hence
\begin{align}
\dot{V} &= \w^T \big( \vect{\tau} - \w \times I\w \big) + K_p\, \sgn(q_{e0})\, \vect{q}_{ev} \cdot \w \nonumber \\
&= \w^T \big( \vect{\tau} + K_p\, \sgn(q_{e0})\, \vect{q}_{ev} \big)
\qquad (\text{since } \w \cdot (\w \times I\w) = 0)
\end{align}
Substituting the control law:
\begin{equation}
\dot{V} = -\w^T K_d\, \w \;\le\; 0 \quad \text{for } K_d > 0
\end{equation}
$\dot{V}$ is only negative \emph{semi}-definite (it vanishes whenever $\w = \vect{0}$), so asymptotic convergence follows from LaSalle's invariance principle: $\w \equiv \vect{0}$ requires $\dot{\w} = \vect{0}$, which forces $\vect{q}_{ev} = \vect{0}$, i.e.\ the target attitude. Note that the gyroscopic term does no work and therefore need not be cancelled for regulation.

\subsection{Gain Selection from Settling Time and Damping Ratio}
For small errors $\vect{q}_{ev} \approx \tfrac{1}{2}\vect{\theta}$, so each principal axis behaves as the second-order system
\begin{equation}
I_i \ddot{\theta}_i + K_d\, \dot{\theta}_i + \frac{K_p}{2}\, \theta_i = 0
\quad\Longrightarrow\quad
\omega_n^2 = \frac{K_p}{2 I_i}, \qquad 2 \zeta \omega_n = \frac{K_d}{I_i}
\end{equation}
Using the 2\,\% settling-time approximation $t_s \approx 4 / (\zeta \omega_n)$, and sizing on the largest principal inertia $I_{\text{ref}} = \max_i I_{ii}$:
\begin{equation}
\omega_n = \frac{4}{\zeta\, t_s},
\qquad
K_p = 2\, I_{\text{ref}}\, \omega_n^2
\end{equation}
\begin{equation}
K_d =
\begin{cases}
2 \zeta \omega_n I_{\text{ref}} \cdot \mathbf{1}_{3 \times 3}, & \text{isotropic option} \\[1ex]
\diag\!\left( 2 \zeta \omega_n \sqrt{I_{\text{ref}}\, I_{ii}} \right), & \text{per-axis option}
\end{cases}
\end{equation}
where $\mathbf{1}_{3\times3}$ is the identity matrix. In the per-axis option the scalar $K_p$ gives axis $i$ a natural frequency $\omega_n \sqrt{I_{\text{ref}}/I_{ii}}$, and the $\sqrt{I_{\text{ref}} I_{ii}}$ scaling of $K_d$ makes the damping ratio equal to $\zeta$ on every axis.

\section{Attitude Tracking: Feedforward plus PD}

\subsection{Rate Error}
The reference supplies $\vect{q}_d(t)$, $\w_d(t)$ and $\dot{\w}_d(t)$, the latter two expressed in the desired frame. They are mapped into the body frame with the DCM of the error quaternion, $C_e = C(\vect{q}_e)$:
\begin{equation}
\w_d^B = C_e\, \w_d,
\qquad
\w_e = \w - \w_d^B
\end{equation}

\subsection{Control Law}
\begin{equation}
\vect{\tau} =
\underbrace{\w \times (I\w)}_{\text{gyroscopic cancellation}}
+ \underbrace{I \left( C_e\, \dot{\w}_d - \w_e \times \w_d^B \right)}_{\text{reference acceleration feedforward}}
- \underbrace{K_p\, \sgn(q_{e0})\, \vect{q}_{ev}}_{\text{attitude feedback}}
- \underbrace{K_d\, \w_e}_{\text{rate-error damping}}
\end{equation}
The feedforward bracket is the body-frame time derivative of $\w_d^B$, because $\dot{C}_e = -[\w_e \times]\, C_e$:
\begin{equation}
\frac{d}{dt}\big( C_e\, \w_d \big) = C_e\, \dot{\w}_d - \w_e \times \w_d^B
\end{equation}
Substituting into Euler's equation gives the closed-loop error dynamics
\begin{equation}
I\, \dot{\w}_e = -K_p\, \sgn(q_{e0})\, \vect{q}_{ev} - K_d\, \w_e
\end{equation}
which has the same structure as the regulation problem, so the tracking error converges with the same Lyapunov argument (with $\w$ replaced by $\w_e$). With a perfect model and matched initial conditions the tracking error is identically zero.

\section{Smooth Reference Generation: Quintic Polynomial Ramp}

A step change in the commanded attitude makes $\w_d$ and $\dot{\w}_d$ impulsive and produces torque spikes. Instead, a rotation of total angle $A$ about a fixed unit axis $\hat{\vect{e}}$ is commanded between $t_s$ and $t_s + T_r$ using a $C^2$-smooth profile. With normalised time $\sigma = (t - t_s)/T_r$:
\begin{align}
s(\sigma) &= 10\sigma^3 - 15\sigma^4 + 6\sigma^5 \\
\dot{s} &= \frac{30\sigma^2 - 60\sigma^3 + 30\sigma^4}{T_r} \\
\ddot{s} &= \frac{60\sigma - 180\sigma^2 + 120\sigma^3}{T_r^2}
\end{align}
with $(s, \dot{s}, \ddot{s}) = (0, 0, 0)$ for $t \le t_s$ and $(1, 0, 0)$ for $t \ge t_s + T_r$. This is the unique fifth-order polynomial satisfying the six boundary conditions
\begin{equation}
s(0) = 0, \quad s(1) = 1, \quad \dot{s}(0) = \dot{s}(1) = 0, \quad \ddot{s}(0) = \ddot{s}(1) = 0
\end{equation}
so angle, rate and acceleration are all continuous. The reference signals are
\begin{equation}
\vect{q}_d(t) = \begin{bmatrix} \cos\!\big( \tfrac{1}{2} A\, s \big) \\ \hat{\vect{e}}\, \sin\!\big( \tfrac{1}{2} A\, s \big) \end{bmatrix},
\qquad
\w_d(t) = \hat{\vect{e}}\, A\, \dot{s},
\qquad
\dot{\w}_d(t) = \hat{\vect{e}}\, A\, \ddot{s}
\end{equation}
For a fixed-axis rotation the components of $\w_d$ are identical in the inertial and desired frames. The peak commanded rate is $\lvert \w_d \rvert_{\max} = 1.875\, A / T_r$, reached at $\sigma = \tfrac{1}{2}$.

\section{Integral Action and Anti-Windup (PID)}

\subsection{Integral State}
A PD law cannot reject a constant disturbance torque: at rest, $\vect{\tau}_{\text{dist}} = K_p\, \sgn(q_{e0})\, \vect{q}_{ev}$ leaves a steady-state attitude error. An integral state $\vect{z} \in \mathbb{R}^3$ on the vector part of the error quaternion is therefore added:
\begin{equation}
\dot{\vect{z}} = \sgn(q_{e0})\, \vect{q}_{ev}
\end{equation}
discretised with forward Euler at the control sample time $\Delta t$, including both anti-windup mechanisms:
\begin{equation}
\vect{z}_{k+1} =
\begin{cases}
\clip\!\big( \vect{z}_k + \sgn(q_{e0})\, \vect{q}_{ev}\, \Delta t,\ -z_{\max},\ +z_{\max} \big), & \text{actuator not saturated} \\[1ex]
\vect{z}_k, & \text{actuator saturated}
\end{cases}
\end{equation}

\begin{itemize}
    \item \textbf{Clamping}: Each component of $\vect{z}$ is limited to $\pm z_{\max}$, bounding the integral torque to $K_i z_{\max}$.
    \item \textbf{Conditional Integration}: The integrator is frozen while the actuator reports saturation, so it does not accumulate error that the actuator cannot act upon.
\end{itemize}

\subsection{PID Control Laws}
The integral term is subtracted from the PD laws of the previous sections ($K_p$, $K_i$ may be per-axis vectors applied element-wise, denoted $\odot$):
\begin{align}
\text{Regulation:} \quad \vect{\tau} &= -K_p \odot \sgn(q_{e0})\, \vect{q}_{ev} - K_d\, \w - K_i \odot \vect{z} \\
\text{Tracking:} \quad \vect{\tau} &= \w \times (I\w) + I \left( C_e\, \dot{\w}_d - \w_e \times \w_d^B \right) - K_p \odot \sgn(q_{e0})\, \vect{q}_{ev} - K_d\, \w_e - K_i \odot \vect{z}
\end{align}
Because the controller now has internal state, it is evaluated exactly once per control sample and \emph{not} inside the RK4 stages.

\subsection{Stability Limit on the Integral Gain}
Linearising one axis ($\vect{q}_{ev} \approx \tfrac{1}{2}\theta$) gives the characteristic polynomial
\begin{equation}
s^3 + \frac{K_d}{I}\, s^2 + \frac{K_p}{2I}\, s + \frac{K_i}{2I} = 0
\end{equation}
and the Routh--Hurwitz criterion ($a_2 a_1 > a_0$) yields the upper bound
\begin{equation}
0 < K_i < \frac{K_p\, K_d}{I}
\end{equation}
(neglecting actuator lag, which reduces the real margin further).

\section{Actuator Model}

\subsection{First-Order Lag}
The actuator is modelled as a first-order system with time constant $T_a$:
\begin{equation}
T_a\, \dot{\vect{\tau}}_{\text{act}} + \vect{\tau}_{\text{act}} = \vect{\tau}_{\text{cmd}}
\end{equation}
Discretising with backward (implicit) Euler, which is unconditionally stable for any $\Delta t$:
\begin{equation}
\vect{\tau}_{\text{act}, k+1} = (1 - \alpha_a)\, \vect{\tau}_{\text{act}, k} + \alpha_a\, \vect{\tau}_{\text{cmd}, k},
\qquad
\alpha_a = \frac{\Delta t}{T_a + \Delta t}
\end{equation}

\subsection{Magnitude Saturation}
The limit is applied to the torque vector norm, so the direction of the torque is preserved:
\begin{equation}
\vect{\tau}_{\text{act}} \leftarrow
\begin{cases}
\vect{\tau}_{\text{act}}, & \lVert \vect{\tau}_{\text{act}} \rVert \le \tau_{\max} \\[1ex]
\tau_{\max}\, \dfrac{\vect{\tau}_{\text{act}}}{\lVert \vect{\tau}_{\text{act}} \rVert}, & \lVert \vect{\tau}_{\text{act}} \rVert > \tau_{\max}
\end{cases}
\end{equation}
The flag $\text{saturated} = \big( \lVert \vect{\tau}_{\text{act}} \rVert > \tau_{\max} \big)$ is passed back to the PID controller on the next sample for conditional integration.

\subsection{Torque Limit from Fin Authority}
For the missile, the available control torque is the aerodynamic moment of a fully deflected fin:
\begin{equation}
\tau_{\max} = \bar{q}\, S\, \bar{c}\, \lvert C_{m\delta} \rvert\, \delta_{\max}
\end{equation}

\section{Missile Model: Atmosphere, Inertia and Aerodynamic Moment}

\subsection{Standard Atmosphere}
With $g_0 = 9.80665\ \text{m/s}^2$, $R = 287.05\ \text{J/(kg\,K)}$, $\gamma = 1.4$ and lapse rate $\lambda = 0.0065\ \text{K/m}$, at altitude $h$:
\begin{equation}
T =
\begin{cases}
288.15 - \lambda h, & h \le 11\,000\ \text{m} \\
216.65\ \text{K}, & h > 11\,000\ \text{m}
\end{cases}
\end{equation}
\begin{equation}
p =
\begin{cases}
101\,325 \left( \dfrac{T}{288.15} \right)^{g_0 / (\lambda R)}, & h \le 11\,000\ \text{m} \\[2ex]
22\,632.06\, \exp\!\left( -\dfrac{g_0\, (h - 11\,000)}{R\, T} \right), & h > 11\,000\ \text{m}
\end{cases}
\end{equation}
\begin{itemize}
    \item \textbf{Air Density} (ideal gas law): $\rho = \dfrac{p}{R\, T}$
    \item \textbf{Speed of Sound}: $a = \sqrt{\gamma R T}$
    \item \textbf{Flight Velocity}: $V = M a$
    \item \textbf{Dynamic Pressure}: $\bar{q} = \tfrac{1}{2} \rho V^2$
\end{itemize}

\subsection{Inertia Matrix}
The missile is treated as a uniform solid cylinder of mass $m$, radius $r$ and length $L$:
\begin{equation}
I_{yy} = I_{zz} = \frac{1}{12} m L^2 + \frac{1}{4} m r^2,
\qquad
I_{xx} = \frac{1}{2} m r^2,
\qquad
I = \diag\!\left( I_{xx},\ I_{yy},\ I_{zz} \right)
\end{equation}

\subsection{Angle of Attack}
For the rotation-only demonstration the velocity vector is assumed fixed along the inertial $x$-axis, so the angle of attack equals the pitch angle of the yaw--pitch--roll sequence:
\begin{equation}
\alpha = \theta = \arcsin\!\big( 2 (q_0 q_2 - q_1 q_3) \big)
\end{equation}

\subsection{Pitch-Channel Aerodynamic Moment}
\begin{equation}
C_m = C_{m\alpha}\, \alpha + C_{m\delta}\, \delta,
\qquad
M_y = \bar{q}\, S\, \bar{c}\, C_m,
\qquad
\vect{M}_{\text{aero}}(\vect{q}) = \begin{bmatrix} 0 \\ M_y \\ 0 \end{bmatrix}
\end{equation}
With $C_{m\alpha} < 0$ the moment is statically stabilising: it opposes any pitch away from $\alpha = 0$ and therefore acts as a disturbance that grows with the commanded pitch angle. The fin deflection $\delta$ is set to zero inside this term, because the control moment is supplied through the actuator torque $\vect{\tau}_{\text{act}}$.

\subsection{Imperfect Model for Feedforward}
The controller holds its own copy of the aero model with every coefficient multiplied by a scale factor $\kappa$:
\begin{equation}
\hat{C}_{m\alpha} = \kappa\, C_{m\alpha},
\qquad
\hat{C}_{m\delta} = \kappa\, C_{m\delta},
\qquad
\hat{\vect{M}}_{\text{aero}} = \kappa\, \vect{M}_{\text{aero}}
\end{equation}
With $\kappa = 0.8$ the controller under-estimates the moment by 20\,\%. After cancellation, the residual disturbance $(1 - \kappa)\, \vect{M}_{\text{aero}}$ remains, and it is this residual that the integral term must remove.

\section{Closed-Loop Missile Pitch Simulation}

\subsection{Per-Axis Gain Scaling}
The roll inertia is roughly $90$ times smaller than the pitch/yaw inertia. A single scalar $K_p$ sized on $I_{\max}$ would make the roll loop faster than the actuator bandwidth and unstable. The reference gains are therefore scaled per axis with $s_i = I_{ii} / I_{\max}$:
\begin{equation}
K_{p,i} = s_i\, K_p^{\text{ref}} = 2\, I_{ii}\, \omega_n^2,
\qquad
K_{d,i} = s_i\, K_d^{\text{ref}} = 2 \zeta \omega_n I_{ii},
\qquad
K_{i,i} = s_i\, K_I
\end{equation}
so every axis has the same natural frequency $\omega_n$ and damping ratio $\zeta$.

\subsection{Commanded Torque}
\begin{equation}
\vect{\tau}_{\text{cmd}} = \vect{\tau}_{\text{PID}} - \hat{\vect{M}}_{\text{aero}}(\vect{q})
\end{equation}
where $\vect{\tau}_{\text{PID}}$ is the tracking PID law and the second term cancels the (imperfectly known) aerodynamic moment.

\subsection{Total Torque on the Plant}
\begin{equation}
\vect{\tau}(t, \vect{x}) = \underbrace{\vect{\tau}_{\text{act}, k}}_{\text{held over the sample}} + \underbrace{\vect{M}_{\text{aero}}\big( \vect{q}(t) \big)}_{\text{re-evaluated in each RK4 stage}}
\end{equation}
The actuator output is a zero-order hold over the control sample, whereas the true aerodynamic moment depends on attitude and is continuous, so it is evaluated inside the RK4 stages.

\subsection{Simulation Algorithm}
For $k = 0, 1, \dots, N$ with $N = \operatorname{round}(t_f / \Delta t)$ and $t_k = k\, \Delta t$:
\begin{enumerate}
    \item Read the plant state $\vect{q}_k$, $\w_k$.
    \item Evaluate the quintic reference: $\vect{q}_d(t_k)$, $\w_d(t_k)$, $\dot{\w}_d(t_k)$.
    \item Compute the error quaternion $\vect{q}_e = \vect{q}_d^{*} \otimes \vect{q}_k$.
    \item Update the integral state (frozen if the actuator saturated on the previous sample) and compute $\vect{\tau}_{\text{PID}}$.
    \item Subtract the modelled aero moment: $\vect{\tau}_{\text{cmd}} = \vect{\tau}_{\text{PID}} - \hat{\vect{M}}_{\text{aero}}(\vect{q}_k)$.
    \item Pass $\vect{\tau}_{\text{cmd}}$ through the actuator lag and saturation to obtain $\vect{\tau}_{\text{act}, k}$.
    \item Log $\theta_{\text{err}}$, $\theta$, $\theta_d$, $\vect{\tau}_{\text{cmd}}$, $\vect{\tau}_{\text{act}}$, $\vect{M}_{\text{aero}}$ and $\vect{z}$.
    \item If $k < N$: advance the plant one RK4 step of size $\Delta t$ with the total torque of the previous subsection.
\end{enumerate}
The simulation is run twice, with $K_I > 0$ (PID) and with $K_I = 0$ (PD), to isolate the effect of integral action on the steady-state error.

\section{Simulation Scenarios and Parameters}

\subsection{Scenario 1: Rest-to-Rest Regulation}
A rigid body is regulated from a $70^\circ$ initial rotation about the body $x$-axis to the identity attitude with the PD law (continuous feedback). The Lyapunov function $V(t)$ is evaluated along the trajectory and plotted on a logarithmic scale to confirm monotonic decay.

\subsection{Scenario 2: Smooth Trajectory Tracking}
The same rigid body tracks a $40^\circ$ quintic ramp about the $z$-axis using the feedforward-plus-PD tracking law, starting from rest at the identity attitude.

\begin{center}
\begin{tabular}{lcc}
\toprule
\textbf{Parameter} & \textbf{Scenario 1} & \textbf{Scenario 2} \\
\midrule
Inertia $I$ [kg\,m$^2$] & $\diag(10, 15, 20)$ & $\diag(10, 15, 20)$ \\
Initial attitude & $70^\circ$ about $x$ & identity \\
Initial rate $\w_0$ & $\vect{0}$ & $\vect{0}$ \\
Target & identity & $40^\circ$ about $z$, $t_s = 2$ s, $T_r = 6$ s \\
Settling time $t_s$ / damping $\zeta$ & $6$ s / $0.8$ & $5$ s / $0.85$ \\
$\omega_n = 4 / (\zeta t_s)$ [rad/s] & $0.833$ & $0.941$ \\
$K_p = 2 I_{\text{ref}}\, \omega_n^2$ & $27.78$ & $35.43$ \\
$K_d = 2 \zeta \omega_n I_{\text{ref}}$ (isotropic) & $26.67$ & $32.00$ \\
Step size $h$ / final time & $0.01$ s / $15$ s & $0.01$ s / $15$ s \\
\bottomrule
\end{tabular}
\end{center}

\subsection{Scenario 3: Missile Pitch Tracking with Aero Moment, Actuator and Integral Action}

\begin{center}
\begin{tabular}{lc}
\toprule
\textbf{Parameter} & \textbf{Value} \\
\midrule
Altitude $h$ / Mach number $M$ & $10\,000$ m / $2.5$ \\
Mass $m$, radius $r$, length $L$ & $180$ kg, $0.113$ m, $2.65$ m \\
Reference area $S$, chord $\bar{c}$ & $0.04$ m$^2$, $0.225$ m \\
$C_{Z\alpha}$, $C_{m\alpha}$, $C_{m\delta}$, $C_{Z\delta}$ [per rad] & $-12$, $-15$, $-20$, $-2$ \\
Pitch command & $8^\circ$ about $y$, $t_s = 1$ s, $T_r = 3$ s \\
Settling time / damping ratio & $1$ s / $0.8$ \\
Integral gain $K_I$ / integral limit $z_{\max}$ & $8000$ / $0.5$ \\
Actuator time constant $T_a$ / fin limit $\delta_{\max}$ & $0.08$ s / $15^\circ$ \\
Aero model scale $\kappa$ & $0.8$ \\
Sample time $\Delta t$ / final time & $0.01$ s / $12$ s \\
\bottomrule
\end{tabular}
\end{center}

Quantities that follow from these parameters through the formulas above:

\begin{center}
\begin{tabular}{lc}
\toprule
\textbf{Derived quantity} & \textbf{Value} \\
\midrule
Temperature $T$ & $223.15$ K \\
Pressure $p$ & $26\,436$ Pa \\
Density $\rho$ & $0.4127$ kg/m$^3$ \\
Speed of sound $a$ / velocity $V$ & $299.46$ m/s / $748.65$ m/s \\
Dynamic pressure $\bar{q}$ & $115\,657$ Pa \\
$I_{yy} = I_{zz}$ / $I_{xx}$ & $105.91$ / $1.149$ kg\,m$^2$ (ratio $\approx 92$) \\
$\omega_n$ & $5.0$ rad/s \\
$K_p^{\text{ref}}$ / $K_d^{\text{ref}}$ & $5295.6$ / $847.3$ \\
Integral stability limit $K_p K_d / I_{yy}$ & $42\,365$ ($K_I$ is $\approx 19\,\%$ of it) \\
Actuator torque limit $\tau_{\max}$ & $5450$ N\,m \\
Aero moment at $\alpha = 8^\circ$ & $-2180$ N\,m \\
\bottomrule
\end{tabular}
\end{center}

\section{Symbol Definitions}

\begin{itemize}
    \item $\vect{q}$, $\vect{q}_d$, $\vect{q}_e$ \textbf{(Quaternions):} Current, desired and error attitude (scalar-first, unit norm).
    \item $\w$, $\w_d$, $\w_e$ \textbf{(Angular Velocities):} Body rate, desired rate (desired frame) and rate error (body frame), in rad/s.
    \item $I$ \textbf{(Inertia Matrix):} Body-frame moment of inertia tensor, kg\,m$^2$.
    \item $\vect{\tau}$, $\vect{\tau}_{\text{cmd}}$, $\vect{\tau}_{\text{act}}$ \textbf{(Torques):} Total, commanded and actuator-delivered body torque, N\,m.
    \item $K_p$, $K_d$, $K_i$ \textbf{(Controller Gains):} Proportional, derivative and integral gains.
    \item $\vect{z}$ \textbf{(Integral State):} Accumulated vector part of the error quaternion.
    \item $c$ \textbf{(Constraint Gain):} Quaternion norm stabilisation gain.
    \item $\zeta$, $\omega_n$, $t_s$ \textbf{(Loop Specifications):} Damping ratio, natural frequency and 2\,\% settling time.
    \item $T_a$, $\tau_{\max}$ \textbf{(Actuator Parameters):} Time constant and torque magnitude limit.
    \item $C_e$ \textbf{(Error DCM):} Direction cosine matrix from the desired frame to the body frame.
    \item $\kappa$ \textbf{(Model Scale):} Ratio of the controller's aero coefficients to the true ones.
\end{itemize}

\end{document}
